\chapter{CAN trong ESP32}

\reportsection{4.1 Tổng quan về giao tiếp CAN trên ESP32}

ESP32 là dòng vi điều khiển SoC được sử dụng rộng rãi trong các hệ thống nhúng nhờ tích
hợp đồng thời năng lực xử lý, kết nối không dây và nhiều ngoại vi truyền thông. Bên
cạnh UART, SPI hay I2C, ESP32 còn cung cấp bộ điều khiển CAN cổ điển do Espressif đặt
tên là TWAI (\textit{Two-Wire Automotive Interface}). Đây là nền tảng giúp ESP32 có thể
tham gia trực tiếp vào một mạng CAN thử nghiệm mà không cần bổ sung thêm bộ điều
khiển CAN rời.

\reportsubsection{4.1.1 Bộ điều khiển TWAI}

TWAI trên ESP32 đảm nhiệm lớp điều khiển giao thức của CAN. Khối này xử lý việc đóng
gói khung dữ liệu, phân xử truy cập bus, tạo CRC, xác nhận ACK, theo dõi bộ đếm lỗi và
quản lý các trạng thái truyền thông. Tuy nhiên, TWAI không tạo trực tiếp tín hiệu vi sai
trên đường truyền, vì vậy vẫn cần thêm một CAN transceiver ở lớp vật lý để đưa tín hiệu
ra hai đường CANH và CANL.

\begin{table}[H]
\centering
\begin{tabular}{|p{4.2cm}|p{9.2cm}|}
\hline
\textbf{Thuộc tính} & \textbf{Mô tả} \\ \hline
Chuẩn hỗ trợ & CAN cổ điển theo ISO 11898-1 \\ \hline
Định dạng khung & Standard Frame với ID 11 bit và Extended Frame với ID 29 bit \\ \hline
Dung lượng dữ liệu & Tối đa 8 byte dữ liệu trong mỗi khung \\ \hline
Chức năng tích hợp & Arbitration, CRC, ACK, bộ đếm lỗi, quản lý bus-off và acceptance filter \\ \hline
Khả năng gán chân & Cho phép ánh xạ TX và RX qua GPIO matrix của ESP32 \\ \hline
Giới hạn & Không hỗ trợ CAN FD, do đó không sử dụng payload 64 byte \\ \hline
\end{tabular}
\caption{Đặc điểm chính của giao tiếp TWAI trên ESP32}
\end{table}



\reportsection{4.2 Kiến trúc một nút CAN dùng ESP32}

Một nút CAN sử dụng ESP32 trong đề tài có thể được xem như sự kết hợp giữa khối xử lý
giao thức ở bên trong vi điều khiển và khối phát thu tín hiệu ở bên ngoài. Việc tách rõ hai
khối này đặc biệt quan trọng khi phân tích phần cứng, vì nhiều lỗi triển khai phát sinh từ
nhầm lẫn giữa CAN controller và CAN transceiver.

\reportsubsection{4.2.1 Vai trò của bộ điều khiển CAN}

TWAI đóng vai trò bộ điều khiển CAN nội bộ. Khối này quyết định định dạng khung truyền,
kiểm tra tính hợp lệ của ID, DLC và cờ điều khiển, đồng thời xử lý các cơ chế cốt lõi của
CAN như tranh chấp bus theo mức ưu tiên ID hay tự động phát hiện lỗi. Ở góc nhìn phần
mềm, đây là nơi firmware cấu hình tốc độ bit, chế độ hoạt động và chính sách lọc khung.

\reportsubsection{4.2.2 Vai trò của bộ thu phát CAN}

CAN transceiver là thành phần thuộc lớp vật lý. Nó nhận mức logic TX/RX từ ESP32 rồi
biến đổi thành cặp tín hiệu vi sai CANH/CANL trên bus. Ở chiều ngược lại, transceiver lấy
tín hiệu từ bus và đưa về mức logic để TWAI có thể giải mã. Nếu thiếu transceiver, ESP32
không thể nối trực tiếp vào mạng CAN thực vì không tạo ra được điện áp và đặc tính chống
nhiễu cần thiết của đường truyền vi sai.

\begin{figure}[H]
\centering
\begin{tikzpicture}[
    block/.style={draw, rectangle, minimum width=3.2cm, minimum height=1.1cm, align=center},
    arrow/.style={->, thick}
]
\node[block] (esp) {ESP32\\TWAI Controller};
\node[block, minimum width=4.2cm, right=2.2cm of esp] (trans) {CAN Transceiver\\ TJA1050};
\node[block, minimum width=3.4cm, right=2.2cm of trans] (bus) {CAN Bus\\CANH -- CANL};

\draw[arrow] (esp) -- node[midway, above, align=center, font=\small, fill=white, inner sep=1pt]{Tín hiệu\\TX/RX mức logic} (trans);
\draw[arrow] (trans) -- node[midway, above, align=center, font=\small, fill=white, inner sep=1pt]{Tín hiệu\\vi sai} (bus);
\end{tikzpicture}
\caption{Sơ đồ khối của một nút CAN dùng ESP32}
\end{figure}

\reportsection{4.3 Kết nối phần cứng của mô hình}

Trong mô hình thực nghiệm hiện tại, firmware mặc định sử dụng GPIO26 làm đường truyền
CAN logic và GPIO27 làm đường nhận CAN logic. Hai chân này được nối với module
transceiver ngoài để hình thành một nút CAN hoàn chỉnh.

\begin{table}[H]
\centering
\begin{tabular}{|>{\raggedright\arraybackslash}p{3.2cm}|>{\raggedright\arraybackslash}p{2.4cm}|>{\raggedright\arraybackslash}p{3.2cm}|>{\raggedright\arraybackslash}p{4.2cm}|}
\hline
\textbf{Nhóm tín hiệu} & \textbf{ESP32} & \textbf{Transceiver} & \textbf{Vai trò} \\ \hline
Truyền dữ liệu & GPIO26 & TXD & Đưa dữ liệu từ TWAI sang bộ thu phát \\ \hline
Nhận dữ liệu & GPIO27 & RXD & Đưa dữ liệu từ bộ thu phát về TWAI \\ \hline
Nguồn nuôi & 3.3 V hoặc 5 V & VCC & Phụ thuộc loại module transceiver sử dụng \\ \hline
Mass chung & GND & GND & Tạo mốc điện áp chung cho toàn nút \\ \hline
Đường bus & -- & CANH, CANL & Kết nối ra đường truyền vi sai của mạng CAN \\ \hline
\end{tabular}
\caption{Ánh xạ chân CAN trong mô hình ESP32 của đề tài}
\end{table}

\reportsection{4.4 Cấu trúc và nguyên lý hoạt động của TWAI}

Khối TWAI trong ESP32 có thể được nhìn nhận như tập hợp của nhiều khối chức năng phối
hợp với nhau để hoàn thành quá trình truyền thông CAN. Cách phân rã này giúp việc mô tả
vai trò của phần cứng rõ ràng hơn và phù hợp với cách trình bày học thuật trong báo cáo.

\reportsubsection{4.4.1 Các khối chức năng chính}

\begin{itemize}[leftmargin=1.5cm, itemsep=0.6em]
    \item Khối điều khiển chế độ hoạt động quyết định TWAI đang ở trạng thái truyền nhận
    bình thường, chỉ nghe hoặc kiểm thử không cần ACK.
    \item Khối định thời bit thiết lập tốc độ truyền và các tham số đồng bộ để mọi nút trên
    bus có cùng nhịp giao tiếp.
    \item Khối truyền chịu trách nhiệm phát khung lên bus, phối hợp với cơ chế arbitration
    và theo dõi việc bản tin có được xác nhận thành công hay không.
    \item Khối nhận lấy tín hiệu từ transceiver, giải mã thành khung CAN rồi chuyển dữ liệu
    cho tầng driver phần mềm.
    \item Khối acceptance filter loại bỏ hoặc cho phép các khung theo cấu hình tiếp nhận.
    \item Khối quản lý lỗi duy trì các bộ đếm lỗi, phát hiện bus-off và hỗ trợ chẩn đoán khi
    đường truyền gặp sự cố.
\end{itemize}

\reportsubsection{4.4.2 Quy trình khởi tạo và trao đổi dữ liệu}

Ở mức tổng quát, một nút ESP32 tham gia mạng CAN thường đi qua các bước sau:

\begin{enumerate}
    \item Chọn cặp chân GPIO dùng cho đường TX và RX của TWAI.
    \item Thiết lập tốc độ bit và cấu hình chế độ hoạt động phù hợp với vai trò của nút.
    \item Đặt chính sách lọc khung để xác định phạm vi bản tin cần tiếp nhận.
    \item Khởi động driver TWAI và đưa transceiver vào trạng thái sẵn sàng làm việc.
    \item Thực hiện truyền hoặc nhận khung thông qua hàng đợi của driver.
    \item Theo dõi trạng thái lỗi để kịp thời xử lý khi bus xuất hiện bất thường.
\end{enumerate}

\begin{figure}[H]
\centering
\begin{tikzpicture}[
    node distance=1.15cm,
    block/.style={draw, rectangle, rounded corners, minimum width=8.0cm, minimum height=0.85cm, align=center},
    setupblock/.style={block, fill=blue!12},
    startblock/.style={block, fill=green!12},
    runblock/.style={block, fill=yellow!15},
    arrow/.style={->, thick}
]
\node[setupblock] (s1) {Chọn chân TX/RX và kết nối với transceiver};
\node[setupblock, below of=s1] (s2) {Thiết lập tốc độ bit, chế độ hoạt động và bộ lọc khung};
\node[startblock, below of=s2] (s3) {Khởi động TWAI và đưa nút vào trạng thái sẵn sàng};
\node[runblock, below of=s3] (s4) {Truyền hoặc nhận bản tin CAN trên bus};
\node[runblock, below of=s4] (s5) {Kiểm tra trạng thái lỗi và cập nhật xử lý tầng trên};

\draw[arrow] (s1) -- (s2);
\draw[arrow] (s2) -- (s3);
\draw[arrow] (s3) -- (s4);
\draw[arrow] (s4) -- (s5);
\end{tikzpicture}
\caption{Quy trình tổng quát khi khởi tạo và vận hành TWAI trên ESP32}
\end{figure}

\reportsection{4.5 Các tham số cấu hình quan trọng}

Để một nút ESP32 giao tiếp ổn định trên bus CAN, một số tham số cấu hình có ý nghĩa quyết
định gồm chế độ hoạt động, tốc độ bit và chính sách lọc khung. Đây cũng là ba nhóm tham
số ảnh hưởng trực tiếp đến việc nút được dùng như ECU mô phỏng, nút tấn công hay nút
giám sát trong hệ thống CyberCAN Defense.

\reportsubsection{4.5.1 Chế độ hoạt động}

\begin{itemize}
    \item \textbf{Normal mode:} nút có thể truyền, nhận và phản hồi ACK như một thành
    phần hoạt động bình thường trên bus.
    \item \textbf{No ACK mode:} hữu ích khi kiểm thử đơn lẻ trong môi trường chưa có nút
    khác xác nhận bản tin, nhưng không phù hợp cho vận hành thực tế.
    \item \textbf{Listen Only mode:} nút chỉ nghe bus, không phát bản tin và không gửi ACK,
    phù hợp với vai trò giám sát hoặc IDS vì không làm thay đổi lưu lượng trên bus.
\end{itemize}

Trong phạm vi đề tài, Normal mode phù hợp với các nút đóng vai trò ECU mô phỏng hoặc
nút phát tấn công, còn Listen Only mode phù hợp với các nút chuyên quan sát và thu thập
log để phân tích.

\reportsubsection{4.5.2 Tốc độ bit và sự đồng bộ trên bus}

Tất cả các nút trên cùng bus CAN phải sử dụng cùng tốc độ bit; nếu không, bản tin sẽ
không được giải mã đúng và mạng nhanh chóng phát sinh lỗi. Một số tốc độ thường gặp
được trình bày trong Bảng 4.3.

\begin{table}[H]
\centering
\begin{tabular}{|p{3cm}|p{9.6cm}|}
\hline
\textbf{Tốc độ bit} & \textbf{Ý nghĩa sử dụng} \\ \hline
125 kbps & Phù hợp với mạng tốc độ thấp hoặc đường truyền dài hơn \\ \hline
250 kbps & Thường dùng trong các mô hình nhúng và công nghiệp quy mô nhỏ \\ \hline
\rowcolor{red!20}
500 kbps & Cân bằng giữa tốc độ và độ ổn định, phù hợp với mô hình thí nghiệm của đề tài \\ \hline
1 Mbps & Dùng cho môi trường cự ly ngắn, yêu cầu băng thông cao hơn \\ \hline
\end{tabular}
\caption{Một số tốc độ bit CAN thường gặp}
\end{table}

Mô hình CyberCAN Defense lựa chọn 500 kbps vì đây là mức cấu hình phổ biến, đủ nhanh
để tái hiện các tình huống tấn công như replay hay flood, đồng thời vẫn ổn định với hệ
thống dây nối ngắn trên bàn thí nghiệm.

\reportsection{4.6 Truyền nhận trên ESP32}

Ở phía truyền, firmware chuẩn bị bản tin bằng cách xác định ID, độ dài dữ liệu, kiểu khung
và nội dung payload trước khi đưa vào hàng đợi của driver. Bộ điều khiển TWAI sau đó
thực hiện phát bản tin lên bus và tự động tham gia cơ chế arbitration nếu cùng lúc có nhiều
nút muốn truyền. Ở phía nhận, các khung hợp lệ được giải mã và chuyển vào hàng đợi nhận
để phần mềm tầng trên đọc ra xử lý. Cách tổ chức này đặc biệt phù hợp với hệ thống có
nhiều tác vụ đồng thời như đề tài hiện tại, nơi việc truyền bản tin, phát hiện tấn công và gửi
log có thể diễn ra song song.

Trong CAN cổ điển, một khung dữ liệu có thể dùng Standard ID 11 bit hoặc Extended ID
29 bit. Việc lựa chọn kiểu ID cần được thống nhất ngay từ giai đoạn thiết kế hệ thống để
tránh sai khác khi các nút giao tiếp chéo với nhau.


